\begin{abstract}
Ce travail présente une étude experimentale de la ré-identification de personnes (Re-ID) à travers sept contributions couvrant l'intégralité de la chaine, du prétraitement à l'indexation d'albums sans supervision. Nous évaluons nos approches sur Market-1501 (32\,668 images, 1,501 identités, 6 caméras) et sur un jeu de données personnel collecté via une pipeline de détection YOLOv8n. Dans un premier temps, nous confrontons des descripteurs classiques de vision par ordinateur (HOG, histogramme de couleur, SIFT) à un modèle profond ResNet-50 entraîné avec une perte combinée identitaire et triplet: l'écart est considérable, avec un Rank-1 passant de 0,68\,\% \`{a} 84,41\,\% en faveur de l'apprentissage profond. Une étude d'ablation sur les stratégies de gel du réseau révele que le fine-tuning complet surpasse nettement le gel total du backbone (85,36\,\% vs 7,07\,\% en Rank-1). La comparaison CNN/Transformer indique que ViT-Base/16 dépasse ResNet-50 en Rank-1 (88,63\,\% vs 84,41\,\%) et en mAP (73,84\,\% vs 64,47\,\%), au coût d'un budget paramétrique 3,45$\times$ plus élevé. L'adaptation de domaine de Market-1501 vers un jeu de données personnel montre une amélioration de la mAP de 78,27\,\% \`{a} 89,93\,\% grâce au transfer learning. Enfin, le clustering DBSCAN sur les embeddings ResNet-50 permet une indexation automatique d'albums avec une pureté de 88,95\,\% sur Market-1501 et 94,80\,\% sur le dataset personnel.
\end{abstract}